\documentclass[12pt]{article}

\usepackage[T1]{fontenc}
\usepackage[utf8]{inputenc}
\usepackage[brazil]{babel}
\usepackage{lmodern}
\usepackage{microtype}
\usepackage[a4paper,margin=2.5cm]{geometry}
\usepackage{setspace}
\usepackage{indentfirst}

\setstretch{1.15}

\title{Por que a Doutrina do Laissez-Faire foi Abandonada? \thanks{Texto Traduzido livremente por Luiz Mario Andrade com o auxílio da ferramenta Chat GPT 5.4}}
\author{Irving Fisher}
\date{4 de janeiro de 1907}

\begin{document}

\maketitle

TALVEZ a mudança mais notável pela qual a opinião econômica passou durante os últimos cinquenta anos tenha sido a mudança das doutrinas extremas de \textit{laissez-faire} dos economistas clássicos para as doutrinas modernas de regulamentação governamental e controle social. E, no entanto, tem havido muito pouca tentativa de explicar por que o \textit{laissez-faire} foi tão geralmente abandonado. Seu abandono tem sido gradual e quase inconsciente, não tanto o resultado de qualquer doutrina abstrata rival, mas o efeito cumulativo da experiência, que em centenas de casos individuais colocou os homens frente a frente com as limitações práticas da política de "deixar fazer". O movimento está nos trazendo rapidamente de volta à antiga visão pela qual a economia foi inicialmente chamada de \textit{economia política}.

O renascimento da atividade governamental nos assuntos econômicos deve-se a causas que são em parte políticas e em parte econômicas. Este artigo trata principalmente das causas econômicas e, portanto, apenas notaremos de passagem os principais aspectos políticos do problema. Uma razão para a extensão do controle governamental da indústria é a força crescente do controle governamental em geral e da confiança popular nele. O \textit{laissez-faire} era uma doutrina natural em uma época em que os governos eram fracos e ineficientes. A mudança de poder trouxe mudança na teoria do poder. Encontramos o seguro obrigatório para trabalhadores no fortemente desenvolvido Império Alemão; a regulamentação das tarifas ferroviárias segue-se ao aumento do poder e à centralização do governo. Pode-se até dizer que muito da moderna regulamentação governamental da indústria resultou da tentativa dos governos de estender seus poderes em legítima defesa. Sentiu-se, por exemplo, que se o governo não controlasse as ferrovias, as ferrovias controlariam o governo. A regulamentação governamental aqui assumiu o aspecto de uma luta pela supremacia. Assim como a Inglaterra sente a necessidade de ter uma marinha igual às marinhas combinadas de várias outras potências europeias, os governos sentem que devem superar os agregados corporativos de riqueza com os quais podem ter de lidar.

Se houvesse espaço, poderíamos discutir a questão de até onde o movimento em direção à interferência governamental pode ser lucrativamente levado. A doutrina do socialismo reside no polo oposto extremo da doutrina do \textit{laissez-faire}, e estamos nos movendo em direção ao socialismo perigosamente rápido. No entanto, existem obstáculos insuperáveis ao sucesso dos projetos socialistas. O poder e a eficiência governamentais são limitados e, quando uma classe da sociedade tenta realmente governar outra, há sempre uma tendência à corrupção, ineficiência, falta de adaptabilidade a novas condições e abuso de poder. O socialismo não pode ser colocado em prática sem oposição e, para se manter, o socialismo deve manter a classe oposta em sujeição. Nominalmente, esta sujeição seria um paternalismo benevolente, mas na história política é a experiência universal que o partido no poder, para se entrincheirar contra ataques, logo usurpa mais poder, emprega métodos indefensáveis e opressivos e tenta estabelecer-se no gozo de privilégios egoístas especiais.

Nosso propósito atual, entretanto, é estudar, não o lado político, mas o econômico do problema. A doutrina do \textit{laissez-faire} é que a interferência governamental, em questões econômicas pelo menos, é desnecessária e prejudicial. Às vezes, acrescenta-se como corolário que o governo não deve apenas deixar os indivíduos em paz, mas que os indivíduos também devem deixar uns aos outros em paz. "Viver e deixar viver" e "Cada um por si" são os lemas deste tipo de individualismo. Os defensores do \textit{laissez-faire} extremo sustentam que uma classe não se justifica em impor seus gostos a outra. Eles dizem: não devemos nos intrometer nos assuntos de nossos vizinhos, mesmo que eles estejam desperdiçando suas vidas no que nos parece gratificações triviais, inúteis ou positivamente prejudiciais. Aqueles que amam a arte, a ciência ou a literatura não têm o direito, dizem-nos, de criticar aqueles que se aborrecem com essas coisas, mas amam lutas de prêmios, cavalos velozes, a alta sociedade ou a vida nababesca.

O raciocínio pelo qual essas doutrinas individualistas foram apoiadas pode ser brevemente declarado em duas proposições: primeiro, cada indivíduo é o melhor juiz do que serve ao seu próprio interesse, e o motivo do autointeresse o leva a assegurar o máximo de bem-estar para si mesmo; e, segundo, como a sociedade é meramente a soma dos indivíduos, o esforço de cada um para assegurar o máximo de bem-estar para si tem como efeito necessário assegurar também o máximo de bem-estar para a sociedade como um todo.

À luz da experiência dos últimos cinquenta anos, não é difícil ver onde cada uma dessas duas proposições está em erro. Primeiro, não é verdade que se possa confiar que cada homem persiga seus próprios interesses. Alguns homens precisam de esclarecimento, devido à ignorância do que constitui seus melhores interesses, e outros precisam de contenção, devido à falta de autocontrole em segui-los. A necessidade de esclarecimento e contenção sempre foi reconhecida no caso das crianças, e um exame das condições reais mostrará que elas se aplicam — muitas vezes com força igual — aos adultos.

A liberdade é certamente indispensável em uma sociedade saudável, mas a liberdade insensivelmente se aproxima da licença. Embora a maioria de nós ainda concorde que as leis suntuárias são mal aconselhadas, há certamente boas razões para manter que o tráfico de bebidas alcoólicas deve ser colocado sob alguma restrição, mesmo que apenas por meio de licença elevada. Não é verdade que o bêbado seja o melhor juiz do que é para o seu próprio bem-estar e o de sua família, e é ainda menos verdade que, mesmo quando ele reconhece plenamente suas falhas, ele terá o autocontrole para agir de acordo com esse conhecimento. Portanto, o problema do álcool torna-se uma questão social e também individual. Novamente, não é verdade que pais ignorantes se justifiquem em impor suas ideias de educação sobre seus filhos; daí o problema do trabalho infantil, em vez de concernir apenas ao indivíduo, como outrora se pensou, tem relações importantes e de longo alcance com a sociedade como um todo. Os mesmos princípios se aplicam à restrição do jogo, do vício, à supressão da literatura indecente, à compulsão sobre os proprietários de terras para tornarem as habitações sanitárias e a muitas outras formas de regulamentação governamental.

Mesmo onde a intervenção governamental é impraticável ou desaconselhável, ainda haverá boas razões para tentar a melhoria das condições através da influência de uma classe sobre outra; daí surgem as agitações sociais e os esforços de uma classe para educar ou instruir outra. Neste princípio baseiam-se os grandes movimentos modernos para o melhoramento humano, exemplificados pela Sociedade para o Estudo e Prevenção da Tuberculose, a Sociedade de Profilaxia Sanitária e Moral, a Federação Cívica Nacional, o Instituto Americano de Serviço Social, o Comitê Nacional do Trabalho Infantil, sociedades de temperança, assentamentos universitários, associações de enfermeiras distritais e outras organizações.

Por estranho que possa parecer para aqueles de nós interessados nesses movimentos hoje, o fato é que, há uma geração, muitos deles teriam sido considerados pela dominante Escola de Manchester não apenas como impraticáveis, mas como desnecessários e possivelmente prejudiciais. Os adeptos desta escola pareciam tratar a diferença entre conhecimento e ignorância como uma mera diferença de opinião, com a qual o governo não tem mais preocupação do que com a diferença de credos religiosos. É certamente verdade que as tentativas dos governos de impor o que é considerado a "verdadeira religião" sobre todo o povo sempre se provaram mal aconselhadas; o reconhecimento disso produziu o sentimento moderno de tolerância religiosa. Mas estamos levando a tolerância longe demais quando nos recusamos a corrigir erros que a ciência demonstra serem falsos. Existem sem dúvida milhões de pessoas hoje que zombariam da ideia de que o cuspir indiscriminado é perigoso para a saúde pública, mas seria tolo permitir que seu preconceito ignorante prevaleça. O bacteriologista sabe o que o ignorante não sabe, e todos os esforços devem ser feitos para transmitir esse conhecimento às massas o mais rápido possível após sua descoberta. Não podemos permitir que nenhum dogma de \textit{laissez-faire} nos impeça de verificar a ignorância suicida.

O mundo consiste em duas classes — os instruídos e os ignorantes — e é essencial para o progresso que os primeiros tenham permissão para dominar os últimos. Mas uma vez que admitimos ser apropriado que as classes instruídas deem instrução às não instruídas, começamos a ver uma vista quase sem limites para o possível melhoramento humano. Em vez de considerar o estado atual da sociedade como normal e desejável porque cada homem naturalmente "busca seus próprios melhores interesses", permitimo-nos julgar cada caso real pelo nosso próprio padrão ideal. Este padrão pode diferir amplamente da média do uso real. Devemos sempre distinguir entre o \textit{ideal} ou \textit{normal}, e o \textit{real} ou \textit{médio}.

A média representa meramente as condições como elas são; o normal representa as condições como elas deveriam ser. No entanto, nada é mais comum do que confundir os dois. De fato, na maioria das tabelas antropométricas ou fisiológicas, a palavra "normal" é usada quase sinonimamente com "média". A altura normal do homem, seu peso normal, sua duração normal de vida, sua dieta normal, sua força, etc., são todos identificados com a média.

Desta forma, toda a questão de uma possível melhoria é evitada. Somos impedidos logo no início de perguntar, por exemplo, se os homens em geral são gordos demais, pois o peso médio da humanidade é assumido como "normal". O absurdo de tal procedimento torna-se aparente assim que consideramos casos em que, por consentimento comum, a média e o normal são mantidos como distintos. Por exemplo, o homem adulto médio certamente não tem dentes normais, pois eles geralmente estão metade apodrecidos; nem cabelo normal, pois ele é geralmente meio calvo; nem postura normal, pois ele é geralmente curvado. A saúde média está abaixo da saúde normal, a moralidade média abaixo da moralidade normal. Na ausência de evidências, não temos o direito de assumir que a média e o normal são idênticos, mesmo quando nos faltam os dados sobre os quais basear uma opinião. Foi apenas recentemente, e em consequência do movimento contra a tuberculose, que especialistas perceberam o quão amplamente diferente é o ar médio que respiramos do ar que é normal para a respiração humana, e que investigações mostraram que a dieta média, na América pelo menos, é anormalmente nitrogenada. Em vista de tais revelações, deveríamos estar de mente aberta o suficiente para aceitar evidências — caso sejam oferecidas — de que a duração média da vida é menos da metade da duração normal, e que a eficiência média é menos da metade da eficiência normal.

Aqueles que habitualmente confundem o normal e a média são impedidos de ver a possibilidade de progresso. Eles assumem a posição, tão anticientífica quanto obstrutiva ao progresso, de que "o que quer que seja, está certo", presuntivamente pelo menos, e marcam cada um que se desvia da média como um excêntrico ou um maníaco. A confusão entre o normal e a média leva, assim, à confusão entre o excêntrico e o pioneiro. Um excêntrico ou um maníaco é propriamente uma pessoa que se desvia do normal e é quase o oposto do pioneiro, que se desvia da média, mas em direção ao normal.

Discrepâncias entre a média e o normal podem aplicar-se — de fato, aplicam-se — ao lado econômico, bem como a outros lados da vida. Mas isso a doutrina do \textit{laissez-faire} negava. O mundo como ele era foi pensado como sendo quase, se não absolutamente, o melhor mundo possível. Um exemplo dessa suposição complacente estava no uso do termo "utilidade" para significar a intensidade do desejo que os homens têm pelas coisas. Tanto quanto sei, o único escritor que tentou distinguir sistematicamente entre os desejos dos homens como eles são e como deveriam ser é Pareto, que para este propósito sugeriu um novo termo — \textit{ofelimidade} — para substituir "utilidade" aplicada aos desejos reais do homem, reservando para o termo "utilidade" seu sentido original do que é intrinsecamente desejável. Assim, para um viciado em ópio, o ópio tem um alto grau de ofelimidade, mas nenhuma utilidade. Os economistas ainda não deram ênfase suficiente à distinção entre a utilidade verdadeira e o que Pareto chama de ofelimidade. Toda uma gama de problemas de melhoramento social é aberta através desta distinção. Os economistas receberam com escárnio as sugestões de reforma de Ruskin. Mas, por mais impraticáveis que fossem suas propostas específicas, seu ponto de vista é certamente mais sensato do que o da maioria dos economistas; pois, como Ruskin apontou, é absurdo considerar como equivalente um milhão de dólares de capital investido na cultura do ópio e um milhão de dólares investidos em escolas.

Mas resta considerar uma segunda falácia no \textit{laissez-faire}. Não apenas é falso que os homens, quando deixados sozinhos, sempre seguirão seus melhores interesses, mas é falso que quando o fazem, eles sempre servirão melhor à sociedade. Para Adam Smith, parecia evidente que o homem servia melhor à sociedade quando servia melhor a si mesmo — embora ele certamente admitisse que a regra tinha exceções no caso de ladrões, assassinos e outros que são obviamente inimigos da sociedade. Mas a extensão em que as "harmonias econômicas" clássicas foram levadas por alguns escritores, embora não incluíssem tais pessoas como ladrões entre os trabalhadores beneficentes, foi, no entanto, surpreendente. A defesa de Herbert Spencer da liberdade de cunhagem privada é bem conhecida, embora qualquer pessoa familiarizada com a "Lei de Gresham" saiba quão quimérica tal instituição seria. Uma sugestão ainda mais surpreendente é a que se diz que Molinari fez em certa época, a saber, que até mesmo a função policial do governo deveria ser deixada para mãos privadas, que corpos policiais deveriam ser simplesmente comitês de vigilância voluntários, algo como as companhias de bombeiros à moda antiga, e que a rivalidade entre essas companhias garantiria um serviço melhor do que o obtido agora através da polícia governamental!

Se pararmos para classificar os efeitos sociais das ações individuais, descobriremos que eles se dividem em três grupos: (1) aquelas ações que beneficiam o próprio indivíduo e não têm efeito sobre os outros; (2) aquelas ações que beneficiam o indivíduo e ao mesmo tempo beneficiam a sociedade; (3) aquelas ações que beneficiam o indivíduo enquanto ao mesmo tempo prejudicam a sociedade. É o terceiro grupo que os doutrinários do \textit{laissez-faire} negligenciaram, e especialmente aquela parte do terceiro grupo em que o prejuízo à sociedade supera o benefício ao indivíduo. Como Huxley disse:

\begin{quote}
Suponhamos, entretanto, para fins de argumentação, que aceitemos a proposição de que as funções do estado podem ser adequadamente resumidas no grande mandamento negativo — "Não permitirás que nenhum homem interfira na liberdade de qualquer outro homem" — sou incapaz de ver que a consequência lógica seja qualquer restrição ao poder do governo, como seus defensores implicam. Se meu vizinho de ao lado opta por ter seus esgotos em tal estado que criem uma atmosfera venenosa, que eu respiro correndo o risco de tifo e difteria, ele restringe minha justa liberdade de viver tanto quanto se andasse por aí com uma pistola ameaçando minha vida; se ele tiver permissão para deixar seus filhos não vacinados, ele poderia muito bem ter permissão para deixar pastilhas de estricnina espalhadas pelo meu caminho; e se ele os criar sem instrução e sem treinamento para ganhar a vida, ele está fazendo o seu melhor para restringir minha liberdade, aumentando o fardo dos impostos para o sustento de prisões e asilos, que eu tenho que pagar.
\end{quote}

Quanto mais alto o estado de civilização, mais completamente as ações de um membro do corpo social influenciam todos os outros, e menos possível é para qualquer homem fazer algo errado sem interferir, mais ou menos, na liberdade de todos os seus concidadãos.

Nos exemplos dados por Huxley, os atos reclamados são prejudiciais não apenas à sociedade, mas ao indivíduo. Mas mesmo quando o ato de um indivíduo é realmente para seu próprio benefício, ele pode não ser para o benefício da sociedade. O paradoxo de que as ações inteligentes de um milhão de indivíduos, cada um tentando melhorar sua própria condição, podem resultar em tornar a condição agregada do milhão pior, é ilustrado ao considerar o efeito da ação individual no caso de um prédio em chamas. Quando um teatro está pegando fogo, milhares de indivíduos frenéticos lutam para sair. No pânico, é sem dúvida do melhor interesse de qualquer indivíduo em particular lutar para passar à frente dos outros; se ele não o fizer, ele está muito mais propenso a ser queimado. E, no entanto, nada é mais certo do que o fato de que a própria intensidade de tais esforços no agregado derrota seus próprios fins. A razão é que o efeito do esforço é principalmente relativo; na medida em que alguém se empurra para frente, ele empurra outros para trás.

Existem inúmeros exemplos de ações que beneficiam o indivíduo, mas prejudicam a sociedade, ou beneficiam uma parte da sociedade, mas prejudicam a sociedade como um todo. Assim, a cidade de Chicago, ao desviar os Grandes Lagos para seu novo sistema de esgoto, tendeu a influenciar o nível desses lagos e, assim, afetar economicamente um grande território, incluindo vários estados da União e também o Canadá. Estima-se que o nível dos lagos possa ser afetado em até seis polegadas.

Uma razão para a interferência federal na irrigação é que o suprimento de água é frequentemente controlado por cidadãos de um estado, enquanto a terra pertence a outro estado ou aos Estados Unidos, e a cooperação entre os dois é difícil de assegurar. A água, nas terras áridas do oeste, é um requisito primordial e, sem ela, as terras não têm valor. De um ponto, o Mt. Union, no Parque Yellowstone, três rios começam — o Missouri, o Columbia e o Colorado — fluindo para o Golfo do México, o Oceano Pacífico e o Golfo da Califórnia, e atravessando um grande número de estados e uma vasta extensão de território. Os interesses mútuos dos proprietários ribeirinhos e daqueles afetados pela irrigação dificilmente poderiam ser ajustados meramente através do jogo dos interesses individuais.

Da mesma forma, o ato de um indivíduo ao destruir florestas influencia o clima e o suprimento de água e, consequentemente, afeta outros indivíduos em partes distantes. Onde os indivíduos na comunidade têm permissão para buscar seus próprios interesses, a destruição de florestas em algumas regiões segue inevitavelmente.

Um efeito semelhante foi visto há alguns anos no caso da disputa das focas entre os Estados Unidos e a Grã-Bretanha. O jogo do motivo individual neste caso tendia à extinção real das focas e só pôde ser refreado pelo acordo mútuo das nações para evitar a caça pelágica.

A ação individual não pode ser confiada para fornecer construção à prova de fogo ou de queima lenta como exigido em uma cidade lotada; pois o indivíduo, embora interessado em proteger-se dos incêndios de seus vizinhos, não está interessado em proteger seus vizinhos de seus próprios incêndios; daí a necessidade e a justificativa para as ordenanças municipais de incêndio. Da mesma forma, o carvão mineral, em cidades como Denver, St. Louis e Pittsburg, constitui um verdadeiro incômodo para toda a cidade; e, no entanto, o proprietário individual de uma fábrica está indubitavelmente seguindo seu próprio melhor interesse ao não substituir o carvão mineral ou usar consumidores de fumaça caros. Tais medidas protetivas redundariam grandemente em benefício da comunidade, mas apenas ligeiramente em seu próprio benefício; daí a necessidade e a justificativa para ordenanças contra a fumaça. A ação individual nunca daria origem a um sistema de parques urbanos, ou mesmo a qualquer sistema útil de ruas. E onde os parques existem, como no caso do Battery Park, em Nova York, há uma tendência constante daqueles que buscam seus interesses individuais de invadir sobre eles. Em Hartford e outras cidades, certos parques desapareceram gradualmente desta forma, para grande dano do público.

Nos casos mencionados, de um conflito entre interesses sociais e individuais, as restrições legais tornam-se necessárias. Mas há muitos exemplos em que, por uma razão ou outra, as restrições legais são impraticáveis. Isso é particularmente verdadeiro em casos onde várias nações estão envolvidas. Não pode haver dúvida, por exemplo, de que os exércitos permanentes e as grandes marinhas são um fardo quase intolerável na Europa, e que sua existência tendeu a aumentar o custo de nosso próprio exército e marinha, a três mil milhas de distância. No entanto, na ausência de qualquer autoridade internacional central ou acordo mútuo para promover o desarmamento, deve-se confessar que é do interesse da Alemanha ou da França, cada uma individualmente, manter seu equipamento militar em um nível comparável ao de seus vizinhos. No entanto, o efeito agregado da competição internacional por poder militar é anular a si mesmo; a vantagem e desvantagem agregadas são puramente relativas. As nações estão em uma corrida louca, cada uma para superar a outra. Sendo o seu objetivo puramente de vantagem relativa, tal vantagem pode ser deslocada de uma para outra, mas não pode acumular-se para todos. Um aumento geral na vantagem relativa é uma contradição em termos, de modo que, no final, os competidores como um todo tiveram apenas seu trabalho por suas dores.

Um exemplo econômico do mesmo caráter nacional internacional, e que recebeu pouquíssima atenção, é encontrado no aumento dos metais monetários. A produção e distribuição de ouro e prata é o efeito da ação individual, cada pessoa buscando seus próprios melhores interesses. No entanto, o efeito agregado sobre esses indivíduos pode ser prejudicial. O prejuízo referido não é o prejuízo imaginário de uma "balança comercial desfavorável", que era o bicho-papão dos mercantilistas, mas exatamente o oposto. Uma nação que aumenta seu estoque de dinheiro é sempre e necessariamente uma perdedora. Este aumento custa à nação ou o trabalho de mineração ou as mercadorias enviadas para fora do país, e por este custo não há retorno algum. Supor que o aumento do dinheiro em si é um retorno valioso é cometer a falácia do inflacionismo. O dinheiro é uma mercadoria muito peculiar. Um aumento geral de outras mercadorias é uma vantagem para a sociedade, mas um aumento geral de dinheiro não é. O inflacionista raciocina que se um governo pode enriquecer uma pessoa imprimindo papel-moeda e concedendo-o a ela, ele só tem que fazer o mesmo por todos para enriquecer a nação. A ilusão do papel-moeda é muito bem compreendida para exigir comentários. É, no entanto, nem sempre percebido que exatamente o mesmo raciocínio se aplica a toda a inflação, mesmo à inflação que a própria natureza cria quando ela libera seus tesouros de ouro enterrados. Os Estados Unidos agora têm \$33 de dinheiro per capita contra \$22 há alguns anos, mas não estamos melhor por conta disso. A menor quantidade de dinheiro é tão útil na troca quanto a maior quantidade. Existem, é claro, males de transição na contração ou expansão da moeda, mas enquanto o nível de preços permanecer constante ou certo, o número absoluto de dólares do meio circulante é uma questão de indiferença. Segue-se que qualquer esforço despendido em aumentar o estoque de dinheiro é esforço desperdiçado, um esforço sem retorno. Este desperdício é um concomitante necessário do individualismo monetário.

Um caso não dissimile, e que agora está causando muita discussão, é o das tarifas ferroviárias. Aqueles que examinaram o funcionamento da competição no transporte ferroviário reconhecem o fato de que esta competição é da variedade chamada competição "predatória", e que nenhuma tarifa estável ou normal para o transporte, sob a qual os capitalistas consentirão em investir na construção ferroviária, pode ocorrer através de tal competição. Aqueles que defendem a competição como uma cura para os males das tarifas ferroviárias não apreciam a mecânica do problema. O efeito da competição é trazer as tarifas para baixo até o custo de operação; ela não deixa provisão para juros sobre o capital investido na empresa. Se o custo de operação é de um centavo por tonelada-milha, enquanto dois centavos são necessários para incluir receita suficiente para pagar juros sobre o custo original, as tarifas sob competição inevitavelmente cairão abaixo do nível de dois centavos para o nível de um centavo. Pois se assumirmos que a tarifa de dois centavos é por um momento a tarifa vigente, é claro que valeria a pena para qualquer concorrente individual reduzir essa tarifa a fim de desviar o tráfego de seus rivais. Mas assim que ele reduz a tarifa, todos os outros devem fazer o mesmo ou perder seu tráfego. Esta competição é meramente autodefesa, e seu efeito final é prejudicar, não beneficiar, todas as ferrovias que nela se engajam. É uma competição predatória. Para que as tarifas possam ser mantidas no nível de dois centavos em vez do nível de um centavo, a competição deve estar ausente, ou deve ser parcial ou imperfeita. No mundo ferroviário real, a competição está usualmente presente em alguns pontos e ausente em outros. A consequência desta mistura de competição e monopólio é que as tarifas serão determinadas diferentemente para alguns pontos do que para outros, e isso constitui o que é chamado de discriminação local. Em um regime onde o monopólio está presente, a discriminação, não apenas deste caráter local, mas a discriminação quanto a pessoas e mercadorias transportadas, é um resultado natural e inevitável. Não é, obviamente, um resultado desejável; mas não é mais indesejável do que a competição predatória, que é o outro chifre do dilema. Esta competição predatória desencoraja o investimento de capital em novas ferrovias, e os transportadores e consumidores devem, no final, sofrer. Este dilema entre os males do monopólio e da competição leva à regulamentação governamental, embora a eficácia deste remédio não seja tudo o que se poderia desejar. Não é nosso propósito discutir a melhor solução de uma questão tão difícil. Estamos meramente interessados em apontar que este problema das tarifas ferroviárias é parcialmente devido à competição predatória e que a competição predatória é mais um exemplo dos efeitos suicidas de seguir cegamente o autointeresse individual.

Inúmeros outros exemplos poderiam ser dados; contentar-nos-emos, entretanto, com um. Como John Rae apontou, existe uma espécie de competição sutil nos gastos privados, devido à rivalidade social — o desejo de distinção através da riqueza. Frequentemente tem sido observado entre clubes sociais de senhoras que começam com entretenimentos simples, que cada anfitriã sucessiva tenta superar sua predecessora no entretenimento oferecido. Começando com chá e bolo, o clube termina com colações elaboradas e caras, até que produz uma drenagem pesada sobre os recursos de seus membros. De maneira precisamente igual, em uma escala maior, impõe-se um pesado fardo sobre todos nós através da rivalidade social dos indivíduos. Se estudarmos a história de Newport ou resorts de moda semelhantes, descobriremos que as corridas sociais resultaram gradualmente em estabelecer um ritmo que apenas os mais ricos podem manter, e que mesmo para eles os gastos representam custo em vez de satisfação. Este custo frequentemente assume a forma de produzir valores fictícios em artigos meramente porque são "exclusivos". Como John Rae diz:\footnote{'Sociological Theory of Capital', por John Rae, ed. por C. W. Mixter, MacMillan, 1905, p. 247.}

\begin{quote}
Um prato de cérebros de rouxinol dificilmente poderia ser um bocado muito delicioso, mas Adam Smith cita de Plínio o preço pago por um único rouxinol como cerca de £66. Segundo Suetônio, nenhuma refeição custava a Vitélio menos de £2.000... Assim, Adam Smith calcula o custo de certas almofadas de um tipo particular usadas para se apoiar à mesa, em £30.000.
\end{quote}

Nem precisamos extrair nossos exemplos da Roma antiga. A "História do Luxo" de Baudrillart mostrará a tendência de produzir o luxo por rivalidade social em todas as eras. Foi apenas recentemente que um americano em Londres deu um jantar que se disse ter custado \$8.000. A mesa foi colocada em uma grande gôndola veneziana situada no meio de um lago artificial, enquanto em uma gôndola menor próxima uma banda estava estacionada.

Muito tem sido dito ultimamente sobre a importância de viver a vida simples, mas até onde eu sei, não houve nenhuma análise para mostrar por que ela não é vivida. Esta análise revelaria que o fracasso em vivê-la é devido a uma espécie de competição predatória inconsciente na sociedade elegante. Quando San Francisco foi destruída por terremoto e incêndio, muitos comentários foram feitos sobre o fato de que muitos não sentiram suas perdas tanto quanto poderia ter sido antecipado. Uma razão para este resultado é sem dúvida encontrada no fato de que as perdas não foram relativas. Se um único indivíduo tivesse se visto subitamente reduzido de um palácio para uma tenda, seu sentimento de perda e desconforto teria sido grande. Ele não poderia mais retornar ao entretenimento social entre seus antigos associados; ele se sentiria "fora disso" e a inveja corroeria seu peito. Mas após a catástrofe de San Francisco, havia pouco lugar para inveja; todos estavam no mesmo barco. Não houve perda relativa, houve apenas a perda absoluta de confortos materiais, e, por estranho que possa parecer para quem não considerou isso, a perda absoluta é a menor das duas.

É difícil superestimar o imposto que é lançado sobre a sociedade através da corrida social. Não estamos conscientes deste peso, porque, como o peso da atmosfera, ele está sempre pressionando sobre nós. O empresário de Nova York compra um chapéu de seda como uma questão de rotina. Ele não pensa em seu custo como um imposto lançado sobre ele pela sociedade. Ele fica satisfeito porque o chapéu preenche uma necessidade, e ele não considera como essa necessidade se originou. É apenas quando o imposto varia por mudança de lugar, assim como quando a pressão atmosférica varia ao subir uma montanha, que ele se dá conta de sua existência. Se ele se mudar para uma cidade menor onde a corrida social é menos intensa e os líderes na corrida são incapazes de estabelecer um ritmo tão alto, ele descobre que o chapéu alto não é mais \textit{de rigueur}. Ele abandona isso e inúmeras outras despesas e sente que está muito melhor. Um cavalheiro recusou recentemente um salário de \$7.000 em Nova York, preferindo \$4.000 em uma cidade menor, sentindo que poderia comprar não mais satisfação real com o primeiro do que com o último. Os \$3.000 extras significavam simplesmente que custaria mais para se manter no nível de seus vizinhos.

O fardo da corrida social é lançado não apenas sobre os ricos, mas sobre todas as classes. Uma modista em New Haven pensou recentemente em evitar competir com estabelecimentos de chapelaria existentes atendendo ao comércio das balconistas. Para sua surpresa, ela descobriu que a tirania da moda era tão forte entre elas quanto entre as ricas. Ela tentou colocar à venda um grande número de chapéus de \$5 e \$6, mas encontrou grande dificuldade em descartá-los, enquanto os poucos chapéus de \$15 e \$16 encontraram uma venda muito pronta. As balconistas queriam que esses chapéus estivessem "na moda". Recentemente, na França, uma família inteira cometeu suicídio porque perderam o capital que consideravam necessário para manter sua posição social.

Muitos argumentos engenhosos foram feitos para justificar o luxo e em alguns deles pode haver verdade. O fato de que o gasto luxuoso pode ser tão prontamente cortado em tempos difíceis fornece uma espécie de amortecedor contra a carência e a fome. As relações do luxo com o crescimento da população merecem um estudo cuidadoso. Mas quaisquer que sejam os benefícios indiretos do luxo, é certo que ele constitui um imposto sobre a sociedade, e um imposto pesado. Parece também verdade que onde o luxo é maior, a civilização decai.

Se houvesse mais espaço, poderíamos discutir remédios para esta corrida social; mas devemos nos contentar em apenas descrever o fenômeno. Ele exemplifica a maneira pela qual a busca individual de cada um pode criar um fardo para todos.

A partir deste e dos outros exemplos que foram revistos, vemos que a mecânica do individualismo não é tão simples quanto os individualistas assumiram. O antigo individualismo requer duas correções: primeiro, o indivíduo pode frequentemente sofrer interferência em seu próprio interesse, devido à sua ignorância ou à sua falta de autocontrole; segundo, mesmo quando um indivíduo pode ser confiado para seguir seus próprios melhores interesses, não se pode assumir que ele servirá assim aos interesses da sociedade. O reconhecimento desses dois fatos é essencial não apenas para o pensamento claro, mas como preliminar para qualquer solução prática dos grandes problemas do melhoramento humano. Estamos sem dúvida hoje em perigo de experimentação socialística excessiva; mas nada pode ser ganho ignorando ou perdoando os males opostos do individualismo. De fato, a ameaça do socialismo pode ser melhor enfrentada se entendermos e reconhecermos os males que ele pretende remediar. O preliminar para o remédio é o diagnóstico, e um diagnóstico preciso nos salvará do erro de ambos os extremos — o extremo, por um lado, de uma overdose de socialismo e o extremo, por outro lado, de omitir toda a medicação que quer que seja.

\bigskip
\begin{flushright}
IRVING FISHER \\
UNIVERSIDADE DE YALE
\end{flushright}

\section*{LIVROS CIENTÍFICOS}

\textit{Sociology and Social Progress: A Handbook for Students of Sociology.} Compilado por THOMAS NIXON CARVER, Ph.D., LL.D., Professor David A. Wells de Economia Política na Universidade de Harvard. Boston, Ginn and Company [1906]. Pp. vi + 810; 8vo. Preço de lista, \$2,75; preço por correio, \$2,95.

Este é um livro oportuno e valioso. Nestes dias em que as questões sociais estão atraindo a atenção de todos, até mesmo dos especialistas científicos, e quando uma massa indigesta de literatura contemporânea está sendo despejada sobre o público, totalmente incapaz de avaliá-la, é da maior importância que as declarações dos mestres do pensamento, ciência e literatura sobre o assunto sejam tornadas acessíveis a todos como guias para o julgamento público. Fazer isso é o propósito deste volume, e mesmo uma enumeração parcial dos autores e obras que foram consultados é suficiente para indicar o valor da compilação. Os mais importantes são: a "Filosofia Positiva" de Comte (condensação inglesa de Harriet Martineau), a "História da Civilização" de Buckle, o "A Origem do Homem" de Darwin, a "Teoria dos Sentimentos Morais" de Adam Smith, o "Physics and Politics" de Bagehot, o "Outlines of Cosmic Philosophy" de Fiske, o "Dados da Ética" de Herbert Spencer, o "Social Evolution" de Kidd, o "Imitation" de Tarde, o "Hereditary Genius" de Galton, o "O Príncipe" de Maquiavel, a "Política" de Aristóteles.

Quando lembramos que cerca de dois terços do espaço são dedicados a estas obras e uma parte justa a outras que ocupam o segundo lugar apenas em relação a estas, podemos bem perdoar a introdução de um certo número de obras menores e até mesmo algumas bastante insignificantes.

As seleções de grandes obras, o que não é uma tarefa fácil, são feitas criteriosamente. Por exemplo, os três temas mais importantes tratados por Buckle, a saber: a influência das leis físicas sobre a sociedade, o papel do desenvolvimento intelectual e a influência da religião, literatura e governo, são apresentados aqui sem abreviações. Os capítulos de Darwin sobre a seleção sexual em relação ao homem são dados na íntegra, com exceção dos detalhes científicos, de modo que fiquem perfeitamente legíveis. E assim por diante com o resto. Aqueles que leram estas obras em sua juventude e retêm apenas uma vaga impressão delas, têm aqui uma oportunidade de refrescar suas mentes com o que há de melhor nelas, e aqueles que nunca as leram podem ganhar deste \textit{digest} uma ideia bastante adequada delas.

Mas o Dr. Carver pretendeu que o livro, como seu título implica, fosse algo mais do que uma mera compilação. Em primeiro lugar, ele forneceu uma introdução a ele, na qual estabelece de forma tão clara quanto já foi feito o verdadeiro escopo e método da sociologia. Seu tratamento é perfeitamente sensato. Ele é um economista da escola moderna que surgiu da recente definição revisada de valor e que traz as grandes ciências da economia e da sociologia em contato simpático uma com a outra. Se ele dá uma ênfase um tanto indevida ao progresso social, ele apenas faz o que outros, incluindo o presente revisor, fizeram antes de terem se dedicado a um estudo sério das condições da ordem social. A doutrina que ele enfatiza especialmente como sua, e que ele havia estabelecido anteriormente, é expressa nestas palavras:

\begin{quote}
"O conflito de interesses é, portanto, o fato fundamental em nossa vida social... Sem ele, não haveria necessidade de justiça, nem de controle social, nem de ciência moral... Se houvesse harmonia de interesses, o governo seria desnecessário."
\end{quote}

Este é um ponto de vista quase idêntico ao de Huxley. Ele vê a sociedade como algo que requer a repressão do autointeresse individual; ele nega o otimismo complacente dos economistas clássicos, de que o autointeresse e o interesse social são necessariamente harmoniosos. Como mostrado pela introdução e as seleções seguindo-a, Carver encontra a causa da desarmonia social na escassez de bens econômicos, e encontra o remédio na adaptação, tanto do homem ao seu ambiente quanto de seu ambiente ao homem. Em resumo, este manual de sociologia é uma análise das causas e remédios para a desarmonia social.

Este ponto de vista prático e objetivo da sociologia é revigorante após a leitura de tanto do que passa por sociologia na América — uma sociologia de terminologia vazia e definições inúteis de "contato social", "consciência de espécie", e assim por diante. Esperamos que este volume seja adotado não apenas como um livro didático em faculdades, mas que encontre seu caminho na biblioteca de todos os homens pensantes.

\bigskip
\begin{flushright}
IRVING FISHER
\end{flushright}

\end{document}